% !PRINTED none
% section-opening page: "Siebenter Abschnitt. Realität der Wurzeln." heading block with §.76;
% no printed folio number was visible on this page in the scan (chapter-opening layout
% replaces the usual running head with the centred section title). Footer furniture
% "Weber, Algebra. I." and signature "16" are excluded under standing policy.

\begin{center}
\gesperrt{Siebenter Abschnitt.}

\gesperrt{Realität der Wurzeln.}

\rule{3cm}{0.4pt}

\S.\,76.

\gesperrt{Allgemeines über Realität von Gleichungswurzeln und über die Discriminante.}
\end{center}

In diesem Abschnitt werden wir uns mit der Frage beschäftigen, wie viele Wurzeln einer algebraischen Gleichung reell sind. Die Co\"efficienten der Gleichung werden dabei als \gesperrt{reelle Zahlen} vorausgesetzt. Wir wollen solche Gleichungen kurz \gesperrt{reelle Gleichungen} nennen und beginnen mit einigen allgemeinen Betrachtungen.

Die reelle Gleichung $n^{\mathrm{ten}}$ Grades
\[
f(x) = 0
\]
hat, wie wir im dritten Abschnitt gesehen haben, $n$ Wurzeln, die entweder reell oder imaginär, d.\,h. von der Form $\xi + i\eta$ sind, mit reellen $\xi$, $\eta$. Die Zahl von $n$ Wurzeln ergiebt sich aber nur dann allgemein, wenn wir unter Umständen eine Wurzel mehrfach zählen, nämlich $(m+1)$fach, wenn mit $f(x)$ zugleich die $m$ ersten Derivirten von $f(x)$ verschwinden. Da nun ein complexer Ausdruck von der Form $X + iY$ nur dann gleich Null ist, wenn die beiden reellen Bestandtheile $X$ und $Y$ einzeln verschwinden, wenn also mit $X+iY$ zugleich $X-iY$ verschwindet, und da ferner bei reellen Co\"efficienten, wenn $f(\xi+i\eta) = X+iY$ ist, sich $f(\xi-i\eta) = X-iY$ ergiebt, so folgt aus $f(\xi+i\eta)=0$, dass zugleich $f(\xi-i\eta)=0$ sein muss, dass also zu jeder imaginären Wurzel $\xi+i\eta$ eine zweite davon verschiedene imaginäre $\xi-i\eta$, d.\,h. die conjugirt imaginäre Wurzel
%CONT>>
